% ============================================================
% SECCION 6 - GESTION Y TRAZABILIDAD (Danela Arteaga - PE5)
% Fragmento para \input en el informe final
% ============================================================
\section{Gestión y trazabilidad}
\label{sec:gestion-trazabilidad}

\subsection{Cierre del hallazgo DR-01}
\label{sec:cierre-dr01}

La re-inspección interna de esta práctica (22-08-2026) registró como hallazgo crítico DR-01 que
la matriz de trazabilidad no contenía filas para los once requisitos funcionales incorporados en
la Entrega 3 (2A): la última fila de tipo RF llegaba hasta RF-16, mientras que el ERS vigente
declara requisitos hasta RF-27
(\texttt{09\_Cierre\_PE5/registro\_reinspeccion\_PE5.csv}). Esta unidad cierra dicho hallazgo.

La matriz de trazabilidad final (\texttt{04\_Trazabilidad/matriz\_trazabilidad.csv}) recorre la
cadena Fuente/Stakeholder $\rightarrow$ RF/RNF $\rightarrow$ Caso de Uso $\rightarrow$ Clase UML
$\rightarrow$ Estado $\rightarrow$ Criterio de aceptación $\rightarrow$ Caso de prueba conceptual
$\rightarrow$ Historia de usuario, y pasa de 42 a 62 filas: las 42 originales (TR-01 a TR-42),
correspondientes a los dieciséis RF y quince RNF+RD migrados o corregidos en entregas anteriores,
más 20 filas nuevas (TR-43 a TR-62) que trazan los once RF (RF-17 a RF-27) y las nueve
características de calidad restantes (RNF-07 a RNF-15) del ERS vigente.

\subsection{Huérfanos y cadenas cerradas}
\label{sec:huerfanos}

\begin{center}
\begin{longtable}{|p{2.6cm}|p{5.5cm}|p{5cm}|}
\hline
\textbf{Identificador} & \textbf{Causa} & \textbf{Acción tomada} \\
\hline
\endfirsthead
\hline
\textbf{Identificador} & \textbf{Causa} & \textbf{Acción tomada} \\
\hline
\endhead
\hline
\endfoot
\hline
\endlastfoot
TR-19 (RF-15) & Mockup marcado como pendiente para un requisito del componente de IA &
Cerrado como DR-05: RF-15 tiene prioridad \emph{Could Have} en MoSCoW; su interfaz queda
justificada y diferida a la siguiente iteración, sin bloquear el cierre de esta entrega \\
\hline
RF-19, RF-24, RF-25, RF-26 & Prioridad \emph{Debe tener} sin historia de usuario en el backlog &
Enlazadas: se redactaron HU-13 a HU-16 con sus criterios de aceptación en formato Gherkin
(adenda a \texttt{04\_Trazabilidad/HU\_criterios\_aceptacion.md}) \\
\hline
RF-17, RF-18, RF-20, RF-21, RF-22, RF-23, RF-27 & Prioridad menor a \emph{Debe tener}, sin
historia de usuario en el backlog & Justificado por escrito: al ser \emph{Debería/Podría tener},
se difiere su historia formal a una iteración futura; se deja el criterio de verificación
detallado en la matriz para no perder el trabajo de especificación ya hecho \\
\hline
RF-23 & Sin clase de dominio asociada en el diagrama de clases (CD-01) & Reportado como hallazgo
abierto de modelado, no de trazabilidad; ver Sección~\ref{sec:cd01-interpretacion} \\
\hline
UC-20 (diagrama) & Caso de uso \emph{Consult Assigned Tasks} sin ningún RF asociado en el ERS
vigente & Reportado para revisión del equipo; se evitó reutilizar el identificador UC-20 al
proponer los nuevos casos de uso, asignando UC-28 al que hubiera correspondido a RF-20 \\
\hline
\end{longtable}
\end{center}

\subsection{Sincronización backlog $\leftrightarrow$ ERS}
\label{sec:sincronizacion}

En dirección Historia $\rightarrow$ RF, las dieciséis historias del backlog (HU-01 a HU-16, tras
esta unidad) tienen cada una su RF de origen: $16/16 = 100\,\%$. En dirección RF
$\rightarrow$ Historia, de los veintisiete RF del ERS, dieciséis tienen historia asociada:
$16/27 \approx 59{,}3\,\%$. Restringido a los RF de prioridad \emph{Debe tener} —el subconjunto
que la guía exige sincronizar por completo—, el resultado es $16/16 = 100\,\%$, frente al
$12/16 = 75\,\%$ con el que cerró la Entrega 3 (2A).

\subsection{Efecto sobre la métrica M4 (trazabilidad)}
\label{sec:efecto-m4}

Como vista complementaria a la métrica M4 de la auditoría estadística
(Sección~\ref{subsec:m1-m4}), el conteo de requisitos con fila vigente en la
matriz antes de esta unidad era de $28/42$ ($66{,}7\,\%$), por debajo de la
referencia de $\geq 90\,\%$. Con las 20 filas nuevas de esta unidad, la matriz
cubre los 42 requisitos de la línea base auditada: $42/42 = 100\,\%$,
cumpliendo la referencia; la incorporación de las filas de los quince
requisitos de IA del Paso 3 queda documentada como acción de mejora. El detalle
de qué columnas de cada fila nueva quedan completas y cuáles
quedan pendientes de un artefacto visual (mockup, diagrama de caso de uso dibujado) se documenta
fila por fila en la columna \texttt{Estado\_de\_la\_traza} de la matriz.
